\documentclass[10pt,a4paper]{amsart}
\usepackage[margin=19mm]{geometry}
\usepackage{amsmath,amssymb,mathtools}
\usepackage[scr=rsfs,cal=euler]{mathalfa}
\usepackage{xfrac}
\usepackage[unicode,hidelinks,bookmarks=false]{hyperref}
\newcommand{\ie}{i.e.\ }
\newcommand{\cf}{cf.\ }
\newcommand{\resp}{respectively\ }
\newcommand{\Subsection}{\subsection}
\providecommand{\nobreakdash}{\nobreak\hbox{-}}
\setlength{\emergencystretch}{2em}
\multlinegap0pt
\usepackage{mathtools}
\usepackage{xcolor}
\newtheorem{theorem}{Theorem}
\newcommand{\ev}[1]{( #1 )} 
\newcommand{\ew}[1]{\left( #1 \right)} 
\title{A general switching method for constructing E-cospectral hypergraphs}
\author{Aida Abiad, Joshua Cooper, Utku Okur}
\date{Source: arXiv:2604.06016v1 (CC BY 4.0)}
\begin{document}
\maketitle

\noindent\emph{Source and licence.} A general switching method for constructing E-cospectral hypergraphs. Version arXiv:2604.06016v1 (CC BY 4.0), distributed by arXiv under the Creative Commons Attribution 4.0 International licence, http://creativecommons.org/licenses/by/4.0/, as recorded in the arXiv licence metadata for this exact version and shown on the abstract page https://arxiv.org/abs/2604.06016v1. Full licence notice: \texttt{licenses/arxiv-2604.06016-CC-BY-4.0.txt}.\par\smallskip

\noindent\textit{Editorial scope.} Counted unit: the article's theorem thm:main\_theorem (source0.tex lines 385--393). The article's complete original proof of it is delivered with it (source0.tex lines 395--455, one \texttt{proof} environment of 61 lines). No mathematics is written, repaired or re-derived: the statement and the proof are the article's own lines, in the article's own notation, and no retained line is substituted or re-flowed.  No further source-local context is needed: every cross-reference of every retained line resolves inside the two delivered spans.  Nothing was omitted from any retained span, so the package carries no omission record.  No retained line contains a citation, so the wrapper carries no bibliography and no bibliography entry.  No retained line contains a figure environment, an image-inclusion command or drawing code, so no image is delivered and the package declares no asset dependency.  Editorial formatting only, in addition: an AMS \texttt{amsart} preamble that restates the article's own declarations and macros used by the retained lines, namely the environments \texttt{theorem} on separate counters, together with the standard packages those lines need.  The retained environments print in this document as Theorem 1 in order of appearance, whereas the article prints the same items with the numbering of its own full document.  The complete native source of the article as distributed in the arXiv e-print for this exact version is bundled unchanged as \texttt{sources/197994/source0.tex}.
\begin{theorem}
\label{thm:main_theorem}
Let $G = \ev{V\ev{G},E\ev{G}}$ be a $k$-uniform hypergraph with $k\geq 2$. Let $C=\{ v_1,\ldots,v_s \}$ be a subset of $V\ev{G}$. Let $D:= V\ev{G}\setminus C$. Assume that for all $r\in \{1,2,\ldots,k\}$ and for all $A \in \binom{D}{k-r} $, there is some hypergraph $L_r^A \in \mathcal{B}^{r}_{R}$ induced on the set $C$ ($E\ev{L_r^A}$ can be empty) such that 
\begin{equation}
\label{eqn:to_refer}
\text{$\left\{ e \cup A \in E\ev{G}: e\in \binom{C}{r} \right\} = E\ev{L_r^A} \odot \{A\}$. }
\end{equation}
To construct a $k$-uniform hypergraph $H$, for each $r$ and each $A \in \binom{D}{k-r}$, remove the edges $E\ev{L_r^A} \odot \{A\}$ and add the edges $ E\ev{t\ev{L_r^A}} \odot \{A\}$. Then, $G$ and $H$ are $E$-cospectral.
\end{theorem}

\begin{proof}
By the construction of $H$, 
\begin{equation}
\label{eqn:assumption_implies_2}
\text{ $\left\{ e \cup A \in E\ev{H}: e\in \binom{C}{r} \right\} = E\ev{t\ev{L_r^A}} \odot \{A\}$}
\end{equation}
for each $r$ and $A \in \binom{D}{k-r}$. By the definition of $t\ev{L_r^A}$, we know that, for each $A$,
\begin{align*}
& R^\top \mathcal{A}_{L_r^A} R = \mathcal{A}_{t\ev{L_r^A}}.
\end{align*}
Given a fixed $A$, let $r:= k-|A|$. We can see that
\begin{align*}
\ev{ R^\top \mathbb{A}^{s}_r R }_{i_1,\ldots,i_k} [ \mathbf{a}\ev{j_1,\ldots,j_r}& = \chi\ev{ \{ j_1, \ldots, j_r \} \in E\ev{L_r^A} }, \forall j_1,\ldots,j_r ]\\
& \hspace{-2cm}= R^\top \mathcal{A}_{L_r^A} R \\
& \hspace{-2cm}= \mathcal{A}_{t\ev{L_r^A}}  \\
& \hspace{-2cm} = \ev{ \mathbb{A}^{s}_r }_{i_1,\ldots,i_k} [ \mathbf{a}\ev{j_1,\ldots,j_r} = \chi\ev{ \{ j_1, \ldots, j_r \} \in E\ev{t\ev{L_r^A}} }, \forall j_1,\ldots,j_r ].
\end{align*}

Recall that $Q = R\oplus I_{n-s}$. To see that $Q^\top \mathcal{A}_G Q = \mathcal{A}_H$, it is enough to show that their values agree for each index set $\{i_1,\ldots,i_k\}$. Let $1\leq i_1 < \ldots < i_k \leq n$ be chosen. Without loss of generality, we may assume that $1\leq i_1 < \ldots < i_r \leq s < i_{r+1} < \ldots < i_k \leq n$, for some $r$. Define $A:= \{ i_{r+1}, \ldots, i_{k}\}$. We know that
\begin{align*}
& (Q^\top \mathbb{A}^n_k Q)_{ i_1,\ldots,i_k } = \sum_{ \{ j_1,\ldots,j_k\} \in \binom{ [n] }{k} } \mathbf{a}\ev{j_1,\ldots,j_k}\mathrm{per}\ew{ Q^{\top}\big\vert_{ \{ i_1,\ldots,i_k\} \times \{ j_1,\ldots,j_k \} } } 
\end{align*}
Since $(Q^{\top})_{ a,b } = \delta_{a,b}$, for any $s<a,b$, where $\delta$ is the Kronecker symbol, it follows that
\begin{equation}
\label{eqn:assumption_4}
(Q^\top \mathbb{A}^n_k Q)_{ i_1,\ldots,i_k } =  \sum_{ \{ j_1,\ldots,j_r\} \in \binom{ [s] }{r} } \mathbf{a}\ev{j_1,\ldots, j_r, i_{r+1},\ldots ,i_k}\mathrm{per}\ew{ R^{\top}\big\vert_{ \{ i_1,\ldots,i_r\} \times \{ j_1,\ldots,j_r \} } } 
\end{equation}
Hence, we obtain
\begin{align*}
&(Q^\top \mathbb{A}^n_k Q)_{ i_1,\ldots,i_k }[ \mathbf{a}\ev{j_1,\ldots,j_k} = \chi\ev{ \{ j_1, \ldots, j_k \} \in E\ev{G} }, \forall j_1,\ldots,j_k ] \\
& =(Q^\top \mathbb{A}^n_k Q)_{ i_1,\ldots,i_k }
[\mathbf{a}\ev{j_1,\ldots,j_r,i_{r+1},\ldots,i_k} = \chi\ev{ \{ j_1,\ldots,j_r,i_{r+1},\ldots,i_k \} \in E\ev{G} }, \forall j_1,\ldots,j_r]
\end{align*}
By assumption (\ref{eqn:to_refer}), we have
\begin{align*}
&(Q^\top \mathbb{A}^n_k Q)_{ i_1,\ldots,i_k }[ \mathbf{a}\ev{j_1,\ldots,j_k} = \chi\ev{ \{ j_1, \ldots, j_k \} \in E\ev{G} }, \forall j_1,\ldots,j_k ] \\
& = (Q^\top \mathbb{A}^n_k Q)_{ i_1,\ldots,i_k } [ \mathbf{a}\ev{j_1,\ldots, j_k} = \chi\ev{ \{ j_1, \ldots, j_r \} \in E\ev{L_r^A} } \delta\ev{j_{r+1},i_{r+1}} \cdots  \delta\ev{j_{k},i_{k}}, \forall j_1,\ldots,j_r  ] 
\end{align*}

On the other hand, 
\begin{align*}
& \ev{ R^\top \mathbb{A}^{s}_r R }_{i_1,\ldots,i_r} = \sum_{ \{ j_1,\ldots,j_r\} \in \binom{ [s] }{r} } \mathbf{a}\ev{j_1,\ldots, j_r}\mathrm{per}\ew{ R^\top\big\vert_{ \{ i_1,\ldots,i_r\} \times \{ j_1,\ldots,j_r \} } } 
\end{align*}
which implies, by (\ref{eqn:assumption_4}) that
\begin{equation}
\label{eqn:assumption_3}
(Q^\top \mathbb{A}^n_k Q)_{ i_1,\ldots,i_k } [ \mathbf{a}\ev{j_1,\ldots, j_r, i_{r+1},\ldots ,i_k} = \mathbf{a}\ev{j_1,\ldots, j_r}, \forall j_1,\ldots,j_r  ] = \ev{ R^\top \mathbb{A}^{s}_r R }_{i_1,\ldots,i_r}
\end{equation}

Then, 
\begin{align*}
& \ev{ Q^\top \mathcal{A}_G Q }_{ i_1,\ldots,i_k }  = (Q^\top \mathbb{A}^n_k Q)_{ i_1,\ldots,i_k } [ \mathbf{a}\ev{j_1,\ldots,j_k} = \chi\ev{ \{ j_1, \ldots, j_k \} \in E\ev{G} }, \forall j_1,\ldots,j_k ] \\
& = (Q^\top \mathbb{A}^n_k Q)_{ i_1,\ldots,i_k } [ \mathbf{a}\ev{j_1,\ldots, j_k} = \chi\ev{ \{ j_1, \ldots, j_r \} \in E\ev{L_r^A} } \delta\ev{j_{r+1},i_{r+1}} \cdots  \delta\ev{j_{k},i_{k}}, \forall j_1,\ldots,j_k ] \\ 
& = \ev{ R^\top \mathbb{A}^{s}_r R }_{i_1,\ldots,i_r} [ \mathbf{a}\ev{j_1,\ldots,j_r} = \chi\ev{ \{ j_1, \ldots, j_r \} \in E\ev{L_r^A} }, \forall j_1,\ldots,j_r ] \\
& \text{ by (\ref{eqn:assumption_3}) } \\
& =  \ev{ \mathbb{A}^{s}_r }_{i_1,\ldots,i_r} [ \mathbf{a}\ev{j_1,\ldots,j_r} = \chi\ev{ \{ j_1, \ldots, j_r \} \in E\ev{t\ev{L_r^A}} }, \forall j_1,\ldots,j_r ]\\ 
& =  \ev{ \mathbb{A}^{n}_k }_{i_1,\ldots,i_k} [ \mathbf{a}\ev{j_1,\ldots,j_k} = \chi\ev{ \{ j_1, \ldots, j_k \} \in E\ev{H} }, \forall j_1,\ldots,j_k ]\\
& \text{ by (\ref{eqn:assumption_implies_2}), in a similar fashion} \\
& = \ev{ \mathcal{A}_H }_{ i_1,\ldots,i_k }.
\qedhere\end{align*}
\end{proof}

\end{document}
